\section{Evidence and experimental design}
\label{sec:design}
Table~\ref{tab:study_design} distinguishes the data used to learn a model from the evidence used to evaluate it. Neither human survey is used for training or selection. Training seeds, repeated API responses and repeated assignments describe different sources of variability and are never counted as additional people.
\begin{table}[pos=htbp]
\caption{Evidence sources and the unit behind each comparison.}\label{tab:study_design}
\small\renewcommand{\arraystretch}{1.15}
\begin{tabularx}{\linewidth}{@{}p{.23\linewidth}p{.28\linewidth}X@{}}\toprule
Question & Independent units & What is compared \\\midrule
Compression fidelity & Six held-out personas; 398 pairs & Matched objectives and specifications against archived Pro targets \\
Response outside training coverage & Thirty new personas; pilot kept separate & Full and delay-withheld Students against a later Go Flash reference \\
Agreement with people & 332 Singapore and 321 Shanghai respondents & Within-person stated responses; later Pro reference on 24 people per city \\
Transmission through supply & 1,000 fixed synthetic travelers & Assignment rules and four supply arms; paired assignment seeds \\\bottomrule
\end{tabularx}
\end{table}


\subsection{Archived targets and matched training}
\label{sec:controlled_design}
The archived benchmark contains disjoint training, validation and test groups of synthetic traveler profiles, here called personas. The test set has 398 response pairs from six personas. Within each training seed, controlled neural models share initialization, architecture, endpoint targets, pair order and optimization budget. The primary choice-model selector minimizes validation macro-source KL, giving each supervision source a role in selection. A response-based selector and a grid of response-loss weights test sensitivity. MNL-S regularization is selected on the same validation criteria. Departure predictors are selected separately by departure MAE, so timing comparisons do not silently change the choice-selection rule. Appendix~\ref{sec:optimization_support} and Supplementary Section~\ref{supp-sec:supp_controlled} document these comparisons.

The original split also withholds a fare--congestion combination while retaining its constituent factors. We audited endpoints, references and linked contrast units to confirm that the held-out combination does not enter training or validation. This is a compositional test. To test a stronger omission, we additionally remove positive-delay endpoints and entire delay-mechanism groups from training and validation, refit feature statistics, and retrain soft-KL and signed-response models. Budgets are matched between methods within each setting. The full and delay-withheld settings have different retained data and update counts, so their difference does not isolate removal of the intervention family as a pure causal factor.

Archived Teacher targets were collected through the official DeepSeek V4 Pro service, according to the authors' collection records and confirmation. The later same-task survey campaign also requested and returned \texttt{deepseek-v4-pro}. Repeated valid vectors are averaged within state. We preserve these targets as the reference for compression; service names alone do not certify an immutable backend across collection periods.

\subsection{A new synthetic-persona reference test}
\label{sec:new_reference_design}
A separate experiment evaluates the frozen Students on new persona profiles, one fixed numerical trip and a common set of travel conditions. Twelve pilot personas are used only for sample-size planning; thirty different personas form the formal sample. Each has a baseline and eleven changes covering delay, fare, access, combinations, rain and PT infeasibility. Three valid API responses are averaged per state. Profiles are excluded from all prior persona sets. The fixed trip and inherited feasible supply profile limit interpretation to these inputs, rather than geographic or population generalization.

This new reference was supplied by OpenCode Go with requested and returned model \texttt{deepseek-v4.1-flash}. It therefore measures agreement with a different API model, not compression error against the original Pro Teacher. The primary contrast, defined before formal collection, is signed-response minus soft-KL error under full training and KL selection. We average the paired differences over contrasts and training seeds within persona, then bootstrap personas. Delay-family and individual-condition comparisons are descriptive, with pointwise intervals. The pilot is excluded from formal estimation. Appendix~\ref{sec:new_protocol} gives the precision rule, acquisition settings and a post-hoc service-metadata diagnostic. No samples were excluded because of that diagnostic.

\subsection{Independent stated-choice evaluation}
\label{sec:human_design}
The Singapore analysis includes 332 residents and the Shanghai model comparison includes 321 respondents. Each completed ten stated-choice tasks. Recruitment through personal contacts and onward sharing produced convenience and snowball samples; the target is agreement within these samples, not citywide demand. Singapore data were collected from 20 August to 5 September 2026 and Shanghai data from 12 to 18 September 2026. Appendix~\ref{sec:supp_inputs} gives eligibility, sample flow and the tasks.

Singapore respondents saw numerical time-and-cost tables. Most Shanghai conditions were qualitative, so unreported durations require numerical reference values for model evaluation. Those assumptions were defined without reference to the human answers. We examine alternative input specifications and a separate network-derived-input diagnostic. All explicit choices are scored, including choices unavailable under the model's ownership rule. The two Shanghai answers indicating no suitable mode are excluded from four-mode scoring. Human-only sensitivity retains eligible respondents outside the model's input mapping.

Within-person contrasts compare stated mode changes with probability changes. Respondent bootstraps retain each person's complete task set and use shared resamples across models. We report nominal 95\% intervals and family-adjusted intervals over the contrasts in each city. Seed variability is separate. An additional same-input API comparison uses 24 respondents per city, with three valid queries per task and no access to their answers. This later Pro reference supports Eq.~\ref{eq:attribution}; it is distinct from both archived supervision and the new Go synthetic sample.

\paragraph{Ethics and consent statement.} The stated-choice questionnaires were administered as anonymous, minimal-risk surveys. No direct personal identifiers were collected. Before beginning the questionnaire, all participants were shown an electronic information and consent statement describing the study purpose, voluntary participation, intended research use of the responses, and confidentiality protections; only respondents who provided consent proceeded to the survey. Formal institutional ethics approval was not obtained. The study procedures were designed to protect participant privacy and to limit data collection to information necessary for the stated research purpose.

\subsection{Perceived conditions and physical transit supply}
\label{sec:execution_design}
The Helsinki experiments hold a synthetic population of 1,000 travelers fixed and use OSM roads and HSL timetables. The original experiment compares assignment rules under an added perceived PT delay while leaving physical supply unchanged. A four-arm extension distinguishes this input manipulation from an actual service reduction: baseline; perceived delay with the original timetable; reduced service with baseline behavioral predictions frozen; and reduced service with predictions recomputed from its level of service. The reduction removes alternating departures within line-and-stop-sequence groups, updating both the route-search index and the MATSim timetable. Delay is not added again to the reduced-service arms.

The extension uses feasibility-constrained assignment, SA-Student and MNL-S with a fixed independent neural timing predictor, and five paired assignment seeds. Its 40 configurations include verified reuse of 12 historical runs and 28 new runs. Responses are paired within person and seed; assignment seeds are averaged within person before person-level bootstrap resampling. This is a single-iteration execution experiment, with no feedback from experienced conditions to a subsequent behavioral decision. Network-derived active-mode speeds are retained, so journey times are diagnostic outputs, not calibrated predictions of the benefits of a service change. Appendix~\ref{sec:execution_support} and Supplementary Section~\ref{supp-sec:supp_execution} provide the stage definitions and additional controls.
